% 第一章：按照本模板的实际使用顺序介绍最常用的 LaTeX 操作。
\chapter{模板快速上手}

\section{本章目的与阅读方式}

本章用于帮助第一次接触 LaTeX 的学生完成模板配置和基本内容修改，
不属于毕业设计（论文）的正式研究内容。建议在开始写作前完整阅读
一遍，并对照源文件实际操作；掌握基本使用方法后，正式提交论文时
可以删除本章及其在 \texttt{main.tex} 中对应的
\verb|% 第一章：按照本模板的实际使用顺序介绍最常用的 LaTeX 操作。
\chapter{模板快速上手}

\section{本章目的与阅读方式}

本章用于帮助第一次接触 LaTeX 的学生完成模板配置和基本内容修改，
不属于毕业设计（论文）的正式研究内容。建议在开始写作前完整阅读
一遍，并对照源文件实际操作；掌握基本使用方法后，正式提交论文时
可以删除本章及其在 \texttt{main.tex} 中对应的
\verb|% 第一章：按照本模板的实际使用顺序介绍最常用的 LaTeX 操作。
\chapter{模板快速上手}

\section{本章目的与阅读方式}

本章用于帮助第一次接触 LaTeX 的学生完成模板配置和基本内容修改，
不属于毕业设计（论文）的正式研究内容。建议在开始写作前完整阅读
一遍，并对照源文件实际操作；掌握基本使用方法后，正式提交论文时
可以删除本章及其在 \texttt{main.tex} 中对应的
\verb|% 第一章：按照本模板的实际使用顺序介绍最常用的 LaTeX 操作。
\chapter{模板快速上手}

\section{本章目的与阅读方式}

本章用于帮助第一次接触 LaTeX 的学生完成模板配置和基本内容修改，
不属于毕业设计（论文）的正式研究内容。建议在开始写作前完整阅读
一遍，并对照源文件实际操作；掌握基本使用方法后，正式提交论文时
可以删除本章及其在 \texttt{main.tex} 中对应的
\verb|\include{chapters/latex-quickstart}| 命令。

LaTeX 的基本工作方式是“在源文件中描述文档结构，再由编译器生成
PDF”。学生只需使用本章介绍的少量命令组织标题、段落、图表和文献，
模板会自动处理学校要求的字体、字号、行距、编号、目录和页码。

\section{从哪里开始修改}

完成论文写作主要涉及三个位置：
\begin{enumerate}
  \item \texttt{main.tex}：填写封面信息、中英文摘要、致谢和附录，
        并安排各章的先后顺序；
  \item \texttt{chapters/}：保存正文各章，每一章对应一个
        \texttt{.tex} 文件；
  \item \texttt{references.bib}：保存参考文献数据。
\end{enumerate}

\texttt{swutthesis.cls} 负责字体、字号、行距、页眉页脚和目录等版式，
正常写作时不修改该文件。模板会自动完成排版，不需要在正文中重复
设置宋体、黑体、字号、行距或页码。

\section{填写论文基本信息}

\texttt{main.tex} 开头集中放置封面信息。花括号中的示例文字应替换为
论文的实际信息：

% 封面命令与 main.tex 保持一致，便于直接对照修改。
\begin{lstlisting}[language={[LaTeX]TeX}]
\thesistitle{论文完整题目}
\coverthesistitle{封面题目第一行\\封面题目第二行}
\thesissubtitle{副标题}

\studentname{姓名}
\studentid{学号}
\college{学院名称}
\major{专业名称}
\supervisor{指导教师}
\coverdate{二〇二六年六月}

\commitmentdate{2026}{6}{20}
\authorizationdate{2026}{6}{20}
\end{lstlisting}

\verb|\thesistitle| 中填写不换行的完整题目，用于页眉和 PDF 信息。
封面需要换成两行时，在 \verb|\coverthesistitle| 中使用一次
\verb|\\|。没有副标题时，删除或注释 \verb|\thesissubtitle| 所在行。
\verb|\commitmentdate{年}{月}{日}| 和
\verb|\authorizationdate{年}{月}{日}| 分别填写诚信承诺书与版权使用
授权声明的日期。如需在打印后手写日期，可将对应命令写为
\verb|\commitmentdate{}{}{}| 或 \verb|\authorizationdate{}{}{}|。

中文摘要写在 \texttt{cnabstract} 环境中，关键词写入
\verb|\cnkeywords{...}|；英文摘要写在 \texttt{enabstract} 环境中，
关键词写入 \verb|\enkeywords{...}|。直接替换模板中的示例内容即可，
摘要标题、缩进、行距和关键词样式均会自动生成。

\section{建立正文各章}

正文文件统一放在 \texttt{chapters/} 目录中。例如，新建
\texttt{chapters/method.tex} 后，可按以下结构开始写作：

\begin{lstlisting}[language={[LaTeX]TeX},float=htbp]
\chapter{研究方法}

\section{总体方案}
这里填写正文。

\subsection{数据来源}
这里填写正文。

\subsubsection{样本筛选}
这里填写正文。
\end{lstlisting}

四个标题命令在 PDF 中依次显示为“第一章”“1.1”“1.1.1”和“（1）”
等层级。标题文字中不输入编号，目录和编号由模板自动生成。

新建的章节文件还需要在 \texttt{main.tex} 正文区域加入：

\begin{lstlisting}[language={[LaTeX]TeX}]
\include{chapters/method}
\end{lstlisting}

\verb|\include| 的排列顺序就是最终的章节顺序。文件名末尾的
\texttt{.tex} 不写入命令。

\section{输入正文和分段}

正文直接输入在相应标题下方。两段文字之间保留一个空行，模板会自动
设置首行缩进和段落行距：

\begin{lstlisting}[language={[LaTeX]TeX}]
这是第一个自然段。源文件中的普通换行不会开始新段落，
PDF 会根据页面宽度自动断行。

这是第二个自然段。两个自然段之间保留一个空行。
\end{lstlisting}

普通正文不使用 \verb|\\| 换段，也不使用连续空格或空行调整位置。
百分号 \verb|%| 后面的内容只作为源文件注释，不会显示在 PDF 中。
正文中需要显示百分号、下划线或连接符号时，分别写成
\verb|\%|、\verb|\_| 和 \verb|\&|。

\section{插入图片、表格和公式}

图片文件放入 \texttt{assets/} 目录，再复制以下代码并替换文件名和
图题：

\begin{lstlisting}[language={[LaTeX]TeX},float=htbp]
\begin{figure}[htbp]
  \centering
  \includegraphics[width=0.72\textwidth]{assets/图片文件名.png}
  \caption{图片标题}
  \label{fig:图片标识}
\end{figure}

如图~\ref{fig:图片标识} 所示，……
\end{lstlisting}

\verb|\label| 中的标识只在源文件中使用，每个标识应保持唯一。
正文通过 \verb|\ref| 引用图号，图片移动后编号会自动更新。

表格和公式的完整标准写法已放在
\texttt{chapters/academic-writing.tex} 中。实际写作时复制其中一个
完整示例，再替换表题、数据或公式内容即可。表格使用
\texttt{table} 和 \texttt{swuttable}，独立编号的公式使用
\texttt{equation}；编号均由模板自动生成。

\section{引用参考文献}

\texttt{references.bib} 已包含图书、期刊论文、学位论文、标准和网络
资料等示例。新增文献时，复制类型最接近的一条记录，修改作者、题名、
年份等字段，并设置一个不重复的引用键。例如，数据库中已有引用键
\texttt{wu2010film}，正文引用方式为：

\begin{lstlisting}[language={[LaTeX]TeX}]
薄膜生长过程具有明显的阶段性特征\cite{wu2010film}。
\end{lstlisting}

文献序号和文末参考文献表会自动生成。正式论文中删除
\texttt{main.tex} 里的 \verb|\nocite{*}| 后，参考文献表只列出正文
实际引用的记录。

按照学校 Word 模板中的写作要求，正式论文参考文献一般不少于 20 篇，
其中外文文献不少于 5 篇，书籍不多于 5 部，教材不列入参考文献。
文末列出的每条文献都应在正文相应位置实际引用，并核对作者、题名、
年份、页码、网址和访问日期等信息。

\section{编译并查看结果}

如果系统已经安装并配置好 \texttt{latexmk}，可在项目根目录运行：

\begin{lstlisting}[language=bash]
latexmk main.tex
\end{lstlisting}

编译完成后查看 \texttt{main.pdf}。修改标题、图表、公式或参考文献后，
仍使用同一条命令重新编译；\texttt{latexmk} 会自动完成目录、编号、
交叉引用和参考文献所需的多轮处理。

部分 MiKTeX Basic 环境没有预装可直接使用的 \texttt{latexmk}。此时
不需要更换模板，可依次执行下面四条命令：

\begin{lstlisting}[language=bash]
xelatex main.tex
biber main
xelatex main.tex
xelatex main.tex
\end{lstlisting}

第一轮 XeLaTeX 生成辅助文件，Biber 处理参考文献，后两轮 XeLaTeX
更新目录、页码和交叉引用。使用 TeXstudio 时，也可以把构建流程配置
为“XeLaTeX → Biber → XeLaTeX → XeLaTeX”。

写作过程中只维护论文内容，不手工输入章节编号、图表编号、页码或
目录点线，也不通过空格和换行模拟版式。模板中的示例可以直接复制，
替换内容后重新编译即可。

\section{进入论文正文示例}

从下一章开始，文档按照“绪论—研究设计与实施—研究结果与规范表达—
讨论与论文组织—总结与展望”的顺序，演示本科毕业设计（论文）的
基本结构，并给出各部分可以参考的写作内容。后续文字用于说明每一
部分应解决的问题，而不是要求学生原样保留；正式写作时应结合本专业
特点、研究任务和指导教师意见替换为自己的论文内容。
| 命令。

LaTeX 的基本工作方式是“在源文件中描述文档结构，再由编译器生成
PDF”。学生只需使用本章介绍的少量命令组织标题、段落、图表和文献，
模板会自动处理学校要求的字体、字号、行距、编号、目录和页码。

\section{从哪里开始修改}

完成论文写作主要涉及三个位置：
\begin{enumerate}
  \item \texttt{main.tex}：填写封面信息、中英文摘要、致谢和附录，
        并安排各章的先后顺序；
  \item \texttt{chapters/}：保存正文各章，每一章对应一个
        \texttt{.tex} 文件；
  \item \texttt{references.bib}：保存参考文献数据。
\end{enumerate}

\texttt{swutthesis.cls} 负责字体、字号、行距、页眉页脚和目录等版式，
正常写作时不修改该文件。模板会自动完成排版，不需要在正文中重复
设置宋体、黑体、字号、行距或页码。

\section{填写论文基本信息}

\texttt{main.tex} 开头集中放置封面信息。花括号中的示例文字应替换为
论文的实际信息：

% 封面命令与 main.tex 保持一致，便于直接对照修改。
\begin{lstlisting}[language={[LaTeX]TeX}]
\thesistitle{论文完整题目}
\coverthesistitle{封面题目第一行\\封面题目第二行}
\thesissubtitle{副标题}

\studentname{姓名}
\studentid{学号}
\college{学院名称}
\major{专业名称}
\supervisor{指导教师}
\coverdate{二〇二六年六月}

\commitmentdate{2026}{6}{20}
\authorizationdate{2026}{6}{20}
\end{lstlisting}

\verb|\thesistitle| 中填写不换行的完整题目，用于页眉和 PDF 信息。
封面需要换成两行时，在 \verb|\coverthesistitle| 中使用一次
\verb|\\|。没有副标题时，删除或注释 \verb|\thesissubtitle| 所在行。
\verb|\commitmentdate{年}{月}{日}| 和
\verb|\authorizationdate{年}{月}{日}| 分别填写诚信承诺书与版权使用
授权声明的日期。如需在打印后手写日期，可将对应命令写为
\verb|\commitmentdate{}{}{}| 或 \verb|\authorizationdate{}{}{}|。

中文摘要写在 \texttt{cnabstract} 环境中，关键词写入
\verb|\cnkeywords{...}|；英文摘要写在 \texttt{enabstract} 环境中，
关键词写入 \verb|\enkeywords{...}|。直接替换模板中的示例内容即可，
摘要标题、缩进、行距和关键词样式均会自动生成。

\section{建立正文各章}

正文文件统一放在 \texttt{chapters/} 目录中。例如，新建
\texttt{chapters/method.tex} 后，可按以下结构开始写作：

\begin{lstlisting}[language={[LaTeX]TeX},float=htbp]
\chapter{研究方法}

\section{总体方案}
这里填写正文。

\subsection{数据来源}
这里填写正文。

\subsubsection{样本筛选}
这里填写正文。
\end{lstlisting}

四个标题命令在 PDF 中依次显示为“第一章”“1.1”“1.1.1”和“（1）”
等层级。标题文字中不输入编号，目录和编号由模板自动生成。

新建的章节文件还需要在 \texttt{main.tex} 正文区域加入：

\begin{lstlisting}[language={[LaTeX]TeX}]
\include{chapters/method}
\end{lstlisting}

\verb|\include| 的排列顺序就是最终的章节顺序。文件名末尾的
\texttt{.tex} 不写入命令。

\section{输入正文和分段}

正文直接输入在相应标题下方。两段文字之间保留一个空行，模板会自动
设置首行缩进和段落行距：

\begin{lstlisting}[language={[LaTeX]TeX}]
这是第一个自然段。源文件中的普通换行不会开始新段落，
PDF 会根据页面宽度自动断行。

这是第二个自然段。两个自然段之间保留一个空行。
\end{lstlisting}

普通正文不使用 \verb|\\| 换段，也不使用连续空格或空行调整位置。
百分号 \verb|%| 后面的内容只作为源文件注释，不会显示在 PDF 中。
正文中需要显示百分号、下划线或连接符号时，分别写成
\verb|\%|、\verb|\_| 和 \verb|\&|。

\section{插入图片、表格和公式}

图片文件放入 \texttt{assets/} 目录，再复制以下代码并替换文件名和
图题：

\begin{lstlisting}[language={[LaTeX]TeX},float=htbp]
\begin{figure}[htbp]
  \centering
  \includegraphics[width=0.72\textwidth]{assets/图片文件名.png}
  \caption{图片标题}
  \label{fig:图片标识}
\end{figure}

如图~\ref{fig:图片标识} 所示，……
\end{lstlisting}

\verb|\label| 中的标识只在源文件中使用，每个标识应保持唯一。
正文通过 \verb|\ref| 引用图号，图片移动后编号会自动更新。

表格和公式的完整标准写法已放在
\texttt{chapters/academic-writing.tex} 中。实际写作时复制其中一个
完整示例，再替换表题、数据或公式内容即可。表格使用
\texttt{table} 和 \texttt{swuttable}，独立编号的公式使用
\texttt{equation}；编号均由模板自动生成。

\section{引用参考文献}

\texttt{references.bib} 已包含图书、期刊论文、学位论文、标准和网络
资料等示例。新增文献时，复制类型最接近的一条记录，修改作者、题名、
年份等字段，并设置一个不重复的引用键。例如，数据库中已有引用键
\texttt{wu2010film}，正文引用方式为：

\begin{lstlisting}[language={[LaTeX]TeX}]
薄膜生长过程具有明显的阶段性特征\cite{wu2010film}。
\end{lstlisting}

文献序号和文末参考文献表会自动生成。正式论文中删除
\texttt{main.tex} 里的 \verb|\nocite{*}| 后，参考文献表只列出正文
实际引用的记录。

按照学校 Word 模板中的写作要求，正式论文参考文献一般不少于 20 篇，
其中外文文献不少于 5 篇，书籍不多于 5 部，教材不列入参考文献。
文末列出的每条文献都应在正文相应位置实际引用，并核对作者、题名、
年份、页码、网址和访问日期等信息。

\section{编译并查看结果}

如果系统已经安装并配置好 \texttt{latexmk}，可在项目根目录运行：

\begin{lstlisting}[language=bash]
latexmk main.tex
\end{lstlisting}

编译完成后查看 \texttt{main.pdf}。修改标题、图表、公式或参考文献后，
仍使用同一条命令重新编译；\texttt{latexmk} 会自动完成目录、编号、
交叉引用和参考文献所需的多轮处理。

部分 MiKTeX Basic 环境没有预装可直接使用的 \texttt{latexmk}。此时
不需要更换模板，可依次执行下面四条命令：

\begin{lstlisting}[language=bash]
xelatex main.tex
biber main
xelatex main.tex
xelatex main.tex
\end{lstlisting}

第一轮 XeLaTeX 生成辅助文件，Biber 处理参考文献，后两轮 XeLaTeX
更新目录、页码和交叉引用。使用 TeXstudio 时，也可以把构建流程配置
为“XeLaTeX → Biber → XeLaTeX → XeLaTeX”。

写作过程中只维护论文内容，不手工输入章节编号、图表编号、页码或
目录点线，也不通过空格和换行模拟版式。模板中的示例可以直接复制，
替换内容后重新编译即可。

\section{进入论文正文示例}

从下一章开始，文档按照“绪论—研究设计与实施—研究结果与规范表达—
讨论与论文组织—总结与展望”的顺序，演示本科毕业设计（论文）的
基本结构，并给出各部分可以参考的写作内容。后续文字用于说明每一
部分应解决的问题，而不是要求学生原样保留；正式写作时应结合本专业
特点、研究任务和指导教师意见替换为自己的论文内容。
| 命令。

LaTeX 的基本工作方式是“在源文件中描述文档结构，再由编译器生成
PDF”。学生只需使用本章介绍的少量命令组织标题、段落、图表和文献，
模板会自动处理学校要求的字体、字号、行距、编号、目录和页码。

\section{从哪里开始修改}

完成论文写作主要涉及三个位置：
\begin{enumerate}
  \item \texttt{main.tex}：填写封面信息、中英文摘要、致谢和附录，
        并安排各章的先后顺序；
  \item \texttt{chapters/}：保存正文各章，每一章对应一个
        \texttt{.tex} 文件；
  \item \texttt{references.bib}：保存参考文献数据。
\end{enumerate}

\texttt{swutthesis.cls} 负责字体、字号、行距、页眉页脚和目录等版式，
正常写作时不修改该文件。模板会自动完成排版，不需要在正文中重复
设置宋体、黑体、字号、行距或页码。

\section{填写论文基本信息}

\texttt{main.tex} 开头集中放置封面信息。花括号中的示例文字应替换为
论文的实际信息：

% 封面命令与 main.tex 保持一致，便于直接对照修改。
\begin{lstlisting}[language={[LaTeX]TeX}]
\thesistitle{论文完整题目}
\coverthesistitle{封面题目第一行\\封面题目第二行}
\thesissubtitle{副标题}

\studentname{姓名}
\studentid{学号}
\college{学院名称}
\major{专业名称}
\supervisor{指导教师}
\coverdate{二〇二六年六月}

\commitmentdate{2026}{6}{20}
\authorizationdate{2026}{6}{20}
\end{lstlisting}

\verb|\thesistitle| 中填写不换行的完整题目，用于页眉和 PDF 信息。
封面需要换成两行时，在 \verb|\coverthesistitle| 中使用一次
\verb|\\|。没有副标题时，删除或注释 \verb|\thesissubtitle| 所在行。
\verb|\commitmentdate{年}{月}{日}| 和
\verb|\authorizationdate{年}{月}{日}| 分别填写诚信承诺书与版权使用
授权声明的日期。如需在打印后手写日期，可将对应命令写为
\verb|\commitmentdate{}{}{}| 或 \verb|\authorizationdate{}{}{}|。

中文摘要写在 \texttt{cnabstract} 环境中，关键词写入
\verb|\cnkeywords{...}|；英文摘要写在 \texttt{enabstract} 环境中，
关键词写入 \verb|\enkeywords{...}|。直接替换模板中的示例内容即可，
摘要标题、缩进、行距和关键词样式均会自动生成。

\section{建立正文各章}

正文文件统一放在 \texttt{chapters/} 目录中。例如，新建
\texttt{chapters/method.tex} 后，可按以下结构开始写作：

\begin{lstlisting}[language={[LaTeX]TeX},float=htbp]
\chapter{研究方法}

\section{总体方案}
这里填写正文。

\subsection{数据来源}
这里填写正文。

\subsubsection{样本筛选}
这里填写正文。
\end{lstlisting}

四个标题命令在 PDF 中依次显示为“第一章”“1.1”“1.1.1”和“（1）”
等层级。标题文字中不输入编号，目录和编号由模板自动生成。

新建的章节文件还需要在 \texttt{main.tex} 正文区域加入：

\begin{lstlisting}[language={[LaTeX]TeX}]
\include{chapters/method}
\end{lstlisting}

\verb|\include| 的排列顺序就是最终的章节顺序。文件名末尾的
\texttt{.tex} 不写入命令。

\section{输入正文和分段}

正文直接输入在相应标题下方。两段文字之间保留一个空行，模板会自动
设置首行缩进和段落行距：

\begin{lstlisting}[language={[LaTeX]TeX}]
这是第一个自然段。源文件中的普通换行不会开始新段落，
PDF 会根据页面宽度自动断行。

这是第二个自然段。两个自然段之间保留一个空行。
\end{lstlisting}

普通正文不使用 \verb|\\| 换段，也不使用连续空格或空行调整位置。
百分号 \verb|%| 后面的内容只作为源文件注释，不会显示在 PDF 中。
正文中需要显示百分号、下划线或连接符号时，分别写成
\verb|\%|、\verb|\_| 和 \verb|\&|。

\section{插入图片、表格和公式}

图片文件放入 \texttt{assets/} 目录，再复制以下代码并替换文件名和
图题：

\begin{lstlisting}[language={[LaTeX]TeX},float=htbp]
\begin{figure}[htbp]
  \centering
  \includegraphics[width=0.72\textwidth]{assets/图片文件名.png}
  \caption{图片标题}
  \label{fig:图片标识}
\end{figure}

如图~\ref{fig:图片标识} 所示，……
\end{lstlisting}

\verb|\label| 中的标识只在源文件中使用，每个标识应保持唯一。
正文通过 \verb|\ref| 引用图号，图片移动后编号会自动更新。

表格和公式的完整标准写法已放在
\texttt{chapters/academic-writing.tex} 中。实际写作时复制其中一个
完整示例，再替换表题、数据或公式内容即可。表格使用
\texttt{table} 和 \texttt{swuttable}，独立编号的公式使用
\texttt{equation}；编号均由模板自动生成。

\section{引用参考文献}

\texttt{references.bib} 已包含图书、期刊论文、学位论文、标准和网络
资料等示例。新增文献时，复制类型最接近的一条记录，修改作者、题名、
年份等字段，并设置一个不重复的引用键。例如，数据库中已有引用键
\texttt{wu2010film}，正文引用方式为：

\begin{lstlisting}[language={[LaTeX]TeX}]
薄膜生长过程具有明显的阶段性特征\cite{wu2010film}。
\end{lstlisting}

文献序号和文末参考文献表会自动生成。正式论文中删除
\texttt{main.tex} 里的 \verb|\nocite{*}| 后，参考文献表只列出正文
实际引用的记录。

按照学校 Word 模板中的写作要求，正式论文参考文献一般不少于 20 篇，
其中外文文献不少于 5 篇，书籍不多于 5 部，教材不列入参考文献。
文末列出的每条文献都应在正文相应位置实际引用，并核对作者、题名、
年份、页码、网址和访问日期等信息。

\section{编译并查看结果}

如果系统已经安装并配置好 \texttt{latexmk}，可在项目根目录运行：

\begin{lstlisting}[language=bash]
latexmk main.tex
\end{lstlisting}

编译完成后查看 \texttt{main.pdf}。修改标题、图表、公式或参考文献后，
仍使用同一条命令重新编译；\texttt{latexmk} 会自动完成目录、编号、
交叉引用和参考文献所需的多轮处理。

部分 MiKTeX Basic 环境没有预装可直接使用的 \texttt{latexmk}。此时
不需要更换模板，可依次执行下面四条命令：

\begin{lstlisting}[language=bash]
xelatex main.tex
biber main
xelatex main.tex
xelatex main.tex
\end{lstlisting}

第一轮 XeLaTeX 生成辅助文件，Biber 处理参考文献，后两轮 XeLaTeX
更新目录、页码和交叉引用。使用 TeXstudio 时，也可以把构建流程配置
为“XeLaTeX → Biber → XeLaTeX → XeLaTeX”。

写作过程中只维护论文内容，不手工输入章节编号、图表编号、页码或
目录点线，也不通过空格和换行模拟版式。模板中的示例可以直接复制，
替换内容后重新编译即可。

\section{进入论文正文示例}

从下一章开始，文档按照“绪论—研究设计与实施—研究结果与规范表达—
讨论与论文组织—总结与展望”的顺序，演示本科毕业设计（论文）的
基本结构，并给出各部分可以参考的写作内容。后续文字用于说明每一
部分应解决的问题，而不是要求学生原样保留；正式写作时应结合本专业
特点、研究任务和指导教师意见替换为自己的论文内容。
| 命令。

LaTeX 的基本工作方式是“在源文件中描述文档结构，再由编译器生成
PDF”。学生只需使用本章介绍的少量命令组织标题、段落、图表和文献，
模板会自动处理学校要求的字体、字号、行距、编号、目录和页码。

\section{从哪里开始修改}

完成论文写作主要涉及三个位置：
\begin{enumerate}
  \item \texttt{main.tex}：填写封面信息、中英文摘要、致谢和附录，
        并安排各章的先后顺序；
  \item \texttt{chapters/}：保存正文各章，每一章对应一个
        \texttt{.tex} 文件；
  \item \texttt{references.bib}：保存参考文献数据。
\end{enumerate}

\texttt{swutthesis.cls} 负责字体、字号、行距、页眉页脚和目录等版式，
正常写作时不修改该文件。模板会自动完成排版，不需要在正文中重复
设置宋体、黑体、字号、行距或页码。

\section{填写论文基本信息}

\texttt{main.tex} 开头集中放置封面信息。花括号中的示例文字应替换为
论文的实际信息：

% 封面命令与 main.tex 保持一致，便于直接对照修改。
\begin{lstlisting}[language={[LaTeX]TeX}]
\thesistitle{论文完整题目}
\coverthesistitle{封面题目第一行\\封面题目第二行}
\thesissubtitle{副标题}

\studentname{姓名}
\studentid{学号}
\college{学院名称}
\major{专业名称}
\supervisor{指导教师}
\coverdate{二〇二六年六月}

\commitmentdate{2026}{6}{20}
\authorizationdate{2026}{6}{20}
\end{lstlisting}

\verb|\thesistitle| 中填写不换行的完整题目，用于页眉和 PDF 信息。
封面需要换成两行时，在 \verb|\coverthesistitle| 中使用一次
\verb|\\|。没有副标题时，删除或注释 \verb|\thesissubtitle| 所在行。
\verb|\commitmentdate{年}{月}{日}| 和
\verb|\authorizationdate{年}{月}{日}| 分别填写诚信承诺书与版权使用
授权声明的日期。如需在打印后手写日期，可将对应命令写为
\verb|\commitmentdate{}{}{}| 或 \verb|\authorizationdate{}{}{}|。

中文摘要写在 \texttt{cnabstract} 环境中，关键词写入
\verb|\cnkeywords{...}|；英文摘要写在 \texttt{enabstract} 环境中，
关键词写入 \verb|\enkeywords{...}|。直接替换模板中的示例内容即可，
摘要标题、缩进、行距和关键词样式均会自动生成。

\section{建立正文各章}

正文文件统一放在 \texttt{chapters/} 目录中。例如，新建
\texttt{chapters/method.tex} 后，可按以下结构开始写作：

\begin{lstlisting}[language={[LaTeX]TeX},float=htbp]
\chapter{研究方法}

\section{总体方案}
这里填写正文。

\subsection{数据来源}
这里填写正文。

\subsubsection{样本筛选}
这里填写正文。
\end{lstlisting}

四个标题命令在 PDF 中依次显示为“第一章”“1.1”“1.1.1”和“（1）”
等层级。标题文字中不输入编号，目录和编号由模板自动生成。

新建的章节文件还需要在 \texttt{main.tex} 正文区域加入：

\begin{lstlisting}[language={[LaTeX]TeX}]
\include{chapters/method}
\end{lstlisting}

\verb|\include| 的排列顺序就是最终的章节顺序。文件名末尾的
\texttt{.tex} 不写入命令。

\section{输入正文和分段}

正文直接输入在相应标题下方。两段文字之间保留一个空行，模板会自动
设置首行缩进和段落行距：

\begin{lstlisting}[language={[LaTeX]TeX}]
这是第一个自然段。源文件中的普通换行不会开始新段落，
PDF 会根据页面宽度自动断行。

这是第二个自然段。两个自然段之间保留一个空行。
\end{lstlisting}

普通正文不使用 \verb|\\| 换段，也不使用连续空格或空行调整位置。
百分号 \verb|%| 后面的内容只作为源文件注释，不会显示在 PDF 中。
正文中需要显示百分号、下划线或连接符号时，分别写成
\verb|\%|、\verb|\_| 和 \verb|\&|。

\section{插入图片、表格和公式}

图片文件放入 \texttt{assets/} 目录，再复制以下代码并替换文件名和
图题：

\begin{lstlisting}[language={[LaTeX]TeX},float=htbp]
\begin{figure}[htbp]
  \centering
  \includegraphics[width=0.72\textwidth]{assets/图片文件名.png}
  \caption{图片标题}
  \label{fig:图片标识}
\end{figure}

如图~\ref{fig:图片标识} 所示，……
\end{lstlisting}

\verb|\label| 中的标识只在源文件中使用，每个标识应保持唯一。
正文通过 \verb|\ref| 引用图号，图片移动后编号会自动更新。

表格和公式的完整标准写法已放在
\texttt{chapters/academic-writing.tex} 中。实际写作时复制其中一个
完整示例，再替换表题、数据或公式内容即可。表格使用
\texttt{table} 和 \texttt{swuttable}，独立编号的公式使用
\texttt{equation}；编号均由模板自动生成。

\section{引用参考文献}

\texttt{references.bib} 已包含图书、期刊论文、学位论文、标准和网络
资料等示例。新增文献时，复制类型最接近的一条记录，修改作者、题名、
年份等字段，并设置一个不重复的引用键。例如，数据库中已有引用键
\texttt{wu2010film}，正文引用方式为：

\begin{lstlisting}[language={[LaTeX]TeX}]
薄膜生长过程具有明显的阶段性特征\cite{wu2010film}。
\end{lstlisting}

文献序号和文末参考文献表会自动生成。正式论文中删除
\texttt{main.tex} 里的 \verb|\nocite{*}| 后，参考文献表只列出正文
实际引用的记录。

按照学校 Word 模板中的写作要求，正式论文参考文献一般不少于 20 篇，
其中外文文献不少于 5 篇，书籍不多于 5 部，教材不列入参考文献。
文末列出的每条文献都应在正文相应位置实际引用，并核对作者、题名、
年份、页码、网址和访问日期等信息。

\section{编译并查看结果}

如果系统已经安装并配置好 \texttt{latexmk}，可在项目根目录运行：

\begin{lstlisting}[language=bash]
latexmk main.tex
\end{lstlisting}

编译完成后查看 \texttt{main.pdf}。修改标题、图表、公式或参考文献后，
仍使用同一条命令重新编译；\texttt{latexmk} 会自动完成目录、编号、
交叉引用和参考文献所需的多轮处理。

部分 MiKTeX Basic 环境没有预装可直接使用的 \texttt{latexmk}。此时
不需要更换模板，可依次执行下面四条命令：

\begin{lstlisting}[language=bash]
xelatex main.tex
biber main
xelatex main.tex
xelatex main.tex
\end{lstlisting}

第一轮 XeLaTeX 生成辅助文件，Biber 处理参考文献，后两轮 XeLaTeX
更新目录、页码和交叉引用。使用 TeXstudio 时，也可以把构建流程配置
为“XeLaTeX → Biber → XeLaTeX → XeLaTeX”。

写作过程中只维护论文内容，不手工输入章节编号、图表编号、页码或
目录点线，也不通过空格和换行模拟版式。模板中的示例可以直接复制，
替换内容后重新编译即可。

\section{进入论文正文示例}

从下一章开始，文档按照“绪论—研究设计与实施—研究结果与规范表达—
讨论与论文组织—总结与展望”的顺序，演示本科毕业设计（论文）的
基本结构，并给出各部分可以参考的写作内容。后续文字用于说明每一
部分应解决的问题，而不是要求学生原样保留；正式写作时应结合本专业
特点、研究任务和指导教师意见替换为自己的论文内容。
